\documentclass[12pt,a4paper]{article}

% ─── Packages ───────────────────────────────────────────────────────────────
\usepackage[utf8]{inputenc}
\usepackage[T1]{fontenc}
\usepackage[french]{babel}
\usepackage{geometry}
\geometry{margin=2.5cm}
\usepackage{graphicx}
\usepackage{amsmath,amssymb}
\usepackage{algorithm}
\usepackage{algpseudocode}
\usepackage{hyperref}
\hypersetup{colorlinks=true, linkcolor=blue, citecolor=blue, urlcolor=blue}
\usepackage{booktabs}
\usepackage{array}
\usepackage{xcolor}
\usepackage{listings}
\usepackage{caption}
\usepackage{subcaption}
\usepackage{cite}
\usepackage{float}
\usepackage{tikz}
\usetikzlibrary{shapes.geometric, arrows, positioning, fit, backgrounds}
\usepackage{pgfplots}
\pgfplotsset{compat=1.18}
\usepackage{enumitem}
\usepackage{fancyhdr}
\usepackage{titlesec}
\usepackage{setspace}

% ─── Style ──────────────────────────────────────────────────────────────────
\setlength{\parskip}{0.5em}
\setlength{\parindent}{1.5em}
\onehalfspacing

\pagestyle{fancy}
\fancyhf{}

\rfoot{\thepage}

\lstset{
  basicstyle=\small\ttfamily,
  breaklines=true,
  frame=single,
  keywordstyle=\color{blue},
  commentstyle=\color{gray},
  stringstyle=\color{red},
  language=Python,
  numbers=left,
  numberstyle=\tiny,
  tabsize=2
}

% ─── Titre ──────────────────────────────────────────────────────────────────
\title{
  \vspace{-1cm}
  \includegraphics[width=3cm]{example-image}\\[0.5cm]
  {\Large Rapport de Projet}\\[0.3cm]
  \rule{\linewidth}{0.4pt}\\[0.3cm]
}



\date{Année universitaire 2025--2026}

% ─── Document ───────────────────────────────────────────────────────────────
\begin{document}

\maketitle
\thispagestyle{empty}

\begin{abstract}
La sécurité des usagers de la route vulnérables (piétons, cyclistes, motocyclistes)
constitue un enjeu majeur dans les réseaux de communication véhiculaires modernes.
Ces usagers peuvent être détectés par des capteurs GPS et des systèmes V2X, permettant
de calculer leur proximité dangereuse avec des véhicules à moteur et de déclencher
des alertes en temps réel. Cependant, les décisions de communication -- transmettre
directement de véhicule à véhicule (V2V) ou transiter par une infrastructure (V2I)
via une unité de bord de route (RSU) -- sont aujourd'hui prises de manière statique,
sans adaptation aux conditions du réseau. Ce rapport propose un mécanisme de sélection
adaptatif du mode de communication fondé sur l'approche \textit{Multi-Armed Bandit} (MAB).
Deux algorithmes sont mis en œuvre et comparés : la politique $\varepsilon$-greedy et
la politique UCB (\textit{Upper Confidence Bound}). Les critères de décision comprennent
le délai de transmission, le taux de perte de paquets et la charge réseau. Les résultats
obtenus par simulation montrent que l'approche MAB permet une sélection dynamique et
efficace du mode de communication, en s'adaptant aux conditions du réseau et au type
d'utilisateur (piéton, motocycliste, véhicule), tout en minimisant la latence de
transmission des alertes de sécurité.

\medskip
\noindent\textbf{Mots-clés :} réseaux véhiculaires, usagers vulnérables de la route, 
multi-armed bandit, $\varepsilon$-greedy, UCB, V2V, V2I, sélection de mode de communication,
ITS-G5, IEEE 802.11p.
\end{abstract}

\newpage
\tableofcontents
\newpage

% ═══════════════════════════════════════════════════════════════════════════
\section{Introduction}
% ═══════════════════════════════════════════════════════════════════════════

\subsection{Contexte et motivation}

Les réseaux véhiculaires (VANETs) et, plus généralement, les communications V2X
(\textit{Vehicle-to-Everything}) représentent un pilier central des systèmes de
transport intelligent (ITS) de nouvelle génération. L'émergence des véhicules connectés
et autonomes a profondément transformé les exigences en matière de communication :
les véhicules ne se contentent plus d'être des entités passives mais deviennent des
nœuds actifs capables d'échanger des données de façon continue avec leur environnement.

Parmi les applications les plus critiques figure la protection des \textbf{usagers de la
route vulnérables} (URV), une catégorie englobant les piétons, les cyclistes, les
motocyclistes et tout utilisateur non protégé par une carrosserie métallique. Selon les
données de l'Union Européenne, ces usagers représentent plus de 47\% des victimes
de la route, malgré leur part relativement faible dans les flux de circulation.
La remontée périodique de données de localisation GPS par les URV, combinée aux
mécanismes d'alerte V2X, offre une fenêtre d'opportunité pour détecter et prévenir
des collisions potentielles.

Cette remontée d'informations peut être traitée à différents niveaux :
\begin{itemize}[leftmargin=*]
  \item \textbf{Au niveau du véhicule}, via des communications directes V2V
        (\textit{Vehicle-to-Vehicle}) standardisées par IEEE 802.11p / ITS-G5 ;
  \item \textbf{Au niveau de l'infrastructure}, via des unités de bord de route
        (RSU -- \textit{RoadSide Units}) qui peuvent agréger, filtrer et redistribuer
        les informations à tous les véhicules dans leur portée (V2I) ;
  \item \textbf{Au niveau du edge cloud}, sur des serveurs déployés en bordure de route,
        capables d'exécuter des algorithmes d'intelligence artificielle pour analyser
        la situation et déclencher des alertes contextuelles.
\end{itemize}

La question se pose alors : \textit{quel mode de communication choisir dans un contexte
donné pour minimiser la latence et garantir la livraison fiable d'une alerte critique ?}

\subsection{Problématique}

Les réseaux véhiculaires se caractérisent par leur haute mobilité, leur densité
variable, et une topologie en perpétuelle évolution. Dans ce contexte, les approches
de sélection statique du mode de communication (V2V ou V2I) s'avèrent insuffisantes.
En effet :
\begin{itemize}[leftmargin=*]
  \item Dans des zones denses (carrefours, intersections), les communications V2V
        peuvent être perturbées par un fort encombrement du canal radio ;
  \item Dans des zones peu denses ou hors de portée des RSUs, les communications V2I
        peuvent ne pas être disponibles ;
  \item La qualité de signal, la charge du réseau et les conditions de mobilité
        varient rapidement et de façon imprévisible.
\end{itemize}

Face à cette incertitude inhérente, nous proposons d'adopter une approche
d'apprentissage par renforcement légère et efficace : le cadre des
\textbf{Multi-Armed Bandits (MAB)}. Cette approche permet d'apprendre en ligne
la meilleure stratégie de communication sans nécessiter de modèle complet de
l'environnement, en équilibrant exploration et exploitation des canaux disponibles.

\subsection{Contributions}

Les contributions de ce projet sont les suivantes :
\begin{enumerate}[leftmargin=*]
  \item Formulation du problème de sélection du mode de communication V2V/V2I
        comme un problème de Multi-Armed Bandit ;
  \item Mise en œuvre et comparaison de deux algorithmes MAB :
        $\varepsilon$-greedy et UCB ;
  \item Conception d'une simulation en Python modélisant des scénarios de
        communication entre URV et véhicules, avec différents profils de priorité ;
  \item Analyse des performances des deux algorithmes selon les critères de
        délai, taux de perte et charge réseau ;
  \item Discussion sur la scalabilité et les perspectives d'intégration dans
        un cadre normatif ETSI ITS-G5.
\end{enumerate}

\subsection{Organisation du rapport}

Le reste de ce rapport est organisé comme suit. La Section~\ref{sec:etat_art}
présente l'état de l'art sur les communications V2X, la protection des URV et
les approches MAB dans les réseaux sans fil. La Section~\ref{sec:scenario}
détaille le scénario d'application et la formulation du problème. La
Section~\ref{sec:methode} décrit en détail l'approche MAB proposée. La
Section~\ref{sec:implementation} présente l'implémentation. Les
Sections~\ref{sec:resultats} et~\ref{sec:discussion} analysent et discutent
les résultats. La Section~\ref{sec:conclusion} conclut le rapport.

% ═══════════════════════════════════════════════════════════════════════════
\section{État de l'Art}
\label{sec:etat_art}
% ═══════════════════════════════════════════════════════════════════════════

\subsection{Communications V2X et standardisation}

Les communications V2X couvrent quatre types d'interactions : V2V (véhicule à
véhicule), V2I (véhicule à infrastructure), V2P (véhicule à piéton) et V2N
(véhicule à réseau). Deux familles technologiques se distinguent actuellement :

\paragraph{DSRC / IEEE 802.11p / ITS-G5.}
Le standard DSRC (\textit{Dedicated Short-Range Communications}) et sa
déclinaison européenne ITS-G5 reposent sur une extension de la norme Wi-Fi
IEEE 802.11a, opérant dans la bande des 5,9~GHz. Ce standard offre une latence
très faible (de l'ordre de quelques millisecondes) et une portée typique de
300 à 1000~m, ce qui le rend particulièrement adapté aux applications de sécurité
routière en temps réel~\cite{etsi_cam}. L'ETSI a standardisé la génération de
messages CAM (\textit{Cooperative Awareness Message}) à une fréquence de 1 à
10~Hz selon la cinématique du véhicule.

\paragraph{C-V2X / 5G NR-V2X.}
Le C-V2X (\textit{Cellular V2X}), introduit par le 3GPP dans les releases 14 et
suivantes, s'appuie sur les réseaux cellulaires LTE et 5G pour offrir une
couverture plus large et un débit plus élevé. Le 5G NR-V2X permet en particulier
une communication en mode \textit{sidelink} (PC5), analogue au V2V, et un mode
infrastructure via le réseau cellulaire.

\subsection{Protection des usagers de la route vulnérables}

La protection des URV fait l'objet d'une attention croissante dans la communauté
de recherche V2X. Les approches se divisent en plusieurs catégories.

\paragraph{Systèmes basés sur GPS et GNSS.}
La remontée périodique de la position GPS d'un URV permet de calculer sa trajectoire
estimée et de détecter une intersection dangereuse avec celle d'un véhicule. Des
travaux comme~\cite{vru_cam} ont étendu le format CAM aux équipements portables des
URV, permettant de transmettre leur localisation à des intervalles réguliers.

\paragraph{Détection par capteurs embarqués.}
Les systèmes ADAS (caméras, LiDAR, radar) permettent de détecter les URV dans
le champ de vision des véhicules. Cependant, ces systèmes sont limités par la
portée des capteurs et les obstructions physiques.

\paragraph{Perception coopérative.}
Les approches de perception coopérative, telles que les CPM (\textit{Collective
Perception Messages}) standardisés par l'ETSI, permettent aux véhicules de partager
leurs détections locales pour construire une représentation collective de
l'environnement. Bouriga et al.~\cite{bouriga2026} ont récemment proposé un cadre
MAB profond pour filtrer les CAMs pertinents au niveau des RSUs, réduisant la
consommation de bande passante de 88\% par rapport à une diffusion naïve.

\paragraph{Alertes inter-RSU.}
Un axe peu exploré consiste à exploiter les communications inter-RSU pour relayer
les informations relatives aux URV d'une zone couverte à une autre. Wang et
al.~\cite{intercoop} ont proposé InterCoop, un mécanisme de perception coopérative
exploitant les interactions spatio-temporelles pour sélectionner les données les
plus pertinentes à partager entre véhicules du réseau.

\subsection{Apprentissage par renforcement dans les réseaux véhiculaires}

L'apprentissage par renforcement (RL) a été largement étudié pour résoudre des
problèmes d'optimisation dans les réseaux véhiculaires dynamiques.

\paragraph{Q-learning et DQN.}
Des travaux comme~\cite{zoghlami2024} ont exploité le Q-learning pour optimiser
les paramètres de génération de messages coopératifs, améliorant la sécurité des
URV tout en minimisant la consommation énergétique. D'autres travaux ont utilisé
des réseaux de neurones profonds (DQN) pour le placement de fonctions réseau
virtuelles dans des contextes au-delà de la 5G~\cite{demendoza2023}.

\paragraph{Multi-Armed Bandits dans les réseaux sans fil.}
Le cadre MAB a été appliqué avec succès à de nombreux problèmes dans les réseaux
sans fil. Dans le contexte V2X, Hashima et al.~\cite{hashima2023} ont proposé un
algorithme basé sur le \textit{Thompson Sampling} pour la sélection opportuniste
de bandes fréquentielles dans des systèmes V2X multi-bandes. Ce travail démontre
l'efficacité des approches MAB pour les décisions en temps réel sans modèle de
l'environnement.

Le MAB contextuel a également été exploré pour la sélection de faisceau dans les
communications V2I mmWave~\cite{cmab_v2i}, où il a permis de doubler le temps de
maintien des connexions par rapport aux approches basées sur le SNR. Boukhalfa
et al.~\cite{boukhalfa2024} ont benchmarké différents algorithmes RL pour le
handover vertical dans des réseaux hybrides véhiculaires.

\paragraph{MAB pour la sélection de mode de communication.}
À notre connaissance, l'application directe des MAB au problème de sélection
binaire V2V/V2I pour la protection des URV reste un sujet peu exploré. Le travail
fondateur de Woodroofe~\cite{woodroofe1979} sur les bandits à un bras avec variable
concomitante, et les extensions multi-bras de~\cite{anantharam2003}, fournissent
les bases théoriques de notre approche. Dans les réseaux V2X, Baccour et
al.~\cite{mab_v2x_risk} ont formulé la gestion du risque d'accident comme un
problème MAB contextuel, montrant que l'agent peut apprendre à évaluer le risque
et adapter son comportement de transmission en conséquence.

\subsection{Synthèse et positionnement}

Le Tableau~\ref{tab:related} synthétise les principaux travaux connexes et positionne
notre contribution. Si des travaux ont appliqué les MAB aux réseaux V2X pour la
sélection de bande, de faisceau ou la gestion du handover, aucun ne traite le
problème spécifique de la sélection adaptative V2V/V2I pour la protection des
usagers vulnérables, en intégrant à la fois les contraintes de priorité des URV
et les métriques réseau multi-critères.

\begin{table}[H]
\centering
\caption{Comparaison des travaux connexes}
\label{tab:related}
\resizebox{\textwidth}{!}{%
\begin{tabular}{@{}lcccccc@{}}
\toprule
\textbf{Travaux} & \textbf{Contexte V2X} & \textbf{MAB} & \textbf{URV} &
\textbf{V2V/V2I} & \textbf{Multi-critères} & \textbf{RSU} \\
\midrule
Bouriga et al.~\cite{bouriga2026}   & \checkmark & \checkmark & --         & V2I seul  & \checkmark & \checkmark \\
Hashima et al.~\cite{hashima2023}   & \checkmark & \checkmark & --         & Multi-bande & \checkmark & \checkmark \\
C-MAB V2I~\cite{cmab_v2i}           & \checkmark & \checkmark & --         & V2I seul  & --         & -- \\
Baccour et al.~\cite{mab_v2x_risk}  & \checkmark & \checkmark & --         & V2V seul  & --         & -- \\
Zoghlami et al.~\cite{zoghlami2024} & \checkmark & --         & \checkmark & V2X       & --         & -- \\
\textbf{Notre travail}              & \checkmark & \checkmark & \checkmark & V2V + V2I & \checkmark & \checkmark \\
\bottomrule
\end{tabular}%
}
\end{table}

% ═══════════════════════════════════════════════════════════════════════════
\section{Scénario et Formulation du Problème}
\label{sec:scenario}
% ═══════════════════════════════════════════════════════════════════════════

\subsection{Description du scénario}

Nous considérons un environnement urbain dans lequel circulent trois types d'entités :
\begin{itemize}[leftmargin=*]
  \item \textbf{Les piétons} (\textit{Pedestrian}) : usagers les plus vulnérables,
        priorité maximale (priorité $= 1$) ;
  \item \textbf{Les motocyclistes} (\textit{Motorcycle}) : priorité intermédiaire
        (priorité $= 2$) ;
  \item \textbf{Les véhicules automobiles} (\textit{Car}) : priorité la plus basse
        parmi les usagers (priorité $= 3$).
\end{itemize}

Chaque entité est modélisée par un objet \texttt{User} dont les attributs incluent
l'identifiant, le protocole de communication, la portée radio, la capacité de
traitement et la priorité. Des infrastructures (RSUs) sont également déployées et
offrent une capacité de traitement supérieure.

La Figure~\ref{fig:scenario} illustre le scénario considéré :

\begin{figure}[H]
\centering
\begin{tikzpicture}[
  vru/.style={circle, draw=orange!80, fill=orange!20, minimum size=0.8cm, font=\small},
  car/.style={rectangle, draw=blue!80, fill=blue!20, minimum size=0.8cm, font=\small},
  rsu/.style={diamond, draw=green!80, fill=green!20, minimum size=1cm, font=\small},
  arrow/.style={->, thick},
  dashed_arrow/.style={->, dashed, thick, gray}
]
  % RSU
  \node[rsu] (rsu) at (0, 0) {RSU};
  % Vehicles
  \node[car] (car1) at (-4, 1)  {Car 1};
  \node[car] (car2) at (4, 1.5) {Car 2};
  \node[car] (car3) at (3, -1.5){Car 3};
  % VRUs
  \node[vru] (ped)  at (-3.5,-1.5) {Piéton};
  \node[vru] (moto) at (0.5, 3)    {Moto};
  % V2V direct
  \draw[arrow, red!70] (ped) -- node[midway, below left, font=\tiny]{V2V} (car1);
  \draw[arrow, red!70] (moto) -- node[midway, right, font=\tiny]{V2V} (car2);
  % V2I paths
  \draw[dashed_arrow, blue!70] (ped) -- node[midway, below, font=\tiny]{V2I} (rsu);
  \draw[dashed_arrow, blue!70] (rsu) -- node[midway, above, font=\tiny]{I2V} (car3);
  % RSU range
  \draw[dashed, green!60, thick] (0,0) circle (2.5cm);
  \node[font=\tiny, green!60] at (1.8, 1.8) {Portée RSU};
\end{tikzpicture}
\caption{Scénario de communication URV-véhicule : mode V2V direct (rouge) et
mode V2I via RSU (bleu)}
\label{fig:scenario}
\end{figure}

Deux modes de communication sont disponibles :
\begin{enumerate}
  \item \textbf{Mode V2V} : communication directe entre l'URV et le(s) véhicule(s)
        à proximité, via IEEE 802.11p. Ce mode est disponible uniquement si le
        destinataire se trouve dans la portée radio de l'émetteur.
  \item \textbf{Mode V2I} : l'URV envoie son message à la RSU la plus proche, qui
        le redistribue à l'ensemble des véhicules dans sa zone de couverture.
        Ce mode offre une couverture plus large mais peut introduire une latence
        supplémentaire en cas de charge réseau élevée.
\end{enumerate}

\subsection{Modèle de communication}

Chaque protocole de communication est caractérisé par les paramètres suivants :
\begin{itemize}[leftmargin=*]
  \item $N$ : charge réseau courante (\textit{network load}) $\in [0, 1]$ ;
  \item $P_l$ : taux de perte de paquets (\textit{packet loss rate}) ;
  \item $P_s$ : taux de succès de transmission (\textit{transmission success rate}) ;
  \item $\tau$ : délai de transmission de base.
\end{itemize}

Pour le mode V2V, le protocole est configuré avec un taux de succès de $P_s^{V2V}
= 0{,}6$ (reflétant les interférences et la mobilité), tandis que le mode V2I
bénéficie d'un $P_s^{V2I} = 0{,}9$ (canal de backhaul plus stable).

Le délai de transmission total pour un message $m$ entre un émetteur $i$ et un
récepteur $j$ à distance $d_{ij}$ est modélisé par :
\begin{equation}
  \tau_{ij}(m) = \tau_0 + \alpha \cdot d_{ij} + \Delta_q(m)
\end{equation}
où $\tau_0$ est le délai de base, $\alpha$ un facteur de propagation, et
$\Delta_q(m)$ le délai de file d'attente accumulé par le message $m$.

\subsection{Formulation du problème de sélection}

\paragraph{Objectif.}
Soit $\mathcal{M} = \{\text{V2V}, \text{V2I}\}$ l'ensemble des modes disponibles.
À chaque instant $t$, un agent de décision sélectionne un mode $a_t \in \mathcal{M}$
pour transmettre une alerte de sécurité. L'objectif est de maximiser la récompense
cumulée sur un horizon $T$ :
\begin{equation}
  \max_{a_1,\ldots,a_T} \sum_{t=1}^T \mathbb{E}\bigl[r_t(a_t)\bigr]
\end{equation}

\paragraph{Récompense.}
La récompense $r_t(a_t)$ est une combinaison pondérée des métriques inverses,
capturant la qualité de la transmission :
\begin{equation}
  r_t(a_t) = \frac{w_\tau}{\tau_t} + \frac{w_l}{P_{l,t}} + \frac{w_N}{N_t}
\end{equation}
où $w_\tau, w_l, w_N$ sont des poids reflétant l'importance relative de chaque
critère (Tableau~\ref{tab:weights}).

\begin{table}[H]
\centering
\caption{Pondération des critères dans la fonction de récompense}
\label{tab:weights}
\begin{tabular}{@{}lcc@{}}
\toprule
\textbf{Critère} & \textbf{Symbole} & \textbf{Poids} \\
\midrule
Délai de transmission & $w_\tau$ & 1{,}0 \\
Taux de perte de paquets & $w_l$ & 1{,}0 \\
Charge réseau & $w_N$ & 0{,}2 \\
\bottomrule
\end{tabular}
\end{table}

Ce problème correspond à un MAB classique à 2 bras (V2V et V2I) avec des
récompenses stochastiques, que nous étendons à un MAB multi-critères en
décomposant chaque bras en trois sous-critères.

% ═══════════════════════════════════════════════════════════════════════════
\section{Méthode Proposée : Multi-Armed Bandit}
\label{sec:methode}
% ═══════════════════════════════════════════════════════════════════════════

\subsection{Rappel sur le cadre MAB}

Le problème du \textit{Multi-Armed Bandit} (MAB), formalisé dès le début du
XX\textsuperscript{e} siècle, modélise les décisions séquentielles sous incertitude.
Un agent dispose de $K$ bras (actions) et, à chaque round $t$, il sélectionne un
bras $a_t \in \{1, \ldots, K\}$ et reçoit une récompense stochastique $r_t \sim
P_{a_t}$ de distribution inconnue $P_k$ avec espérance $\mu_k$.

L'objectif est de maximiser la récompense cumulée, ou de manière équivalente de
minimiser le \textit{regret cumulatif} :
\begin{equation}
  \mathcal{R}(T) = T \cdot \mu^* - \sum_{t=1}^T \mathbb{E}[r_t]
  \quad \text{avec} \quad \mu^* = \max_k \mu_k
\end{equation}

\subsection{Algorithme $\varepsilon$-Greedy}

L'algorithme $\varepsilon$-greedy est l'une des solutions les plus simples et
les plus robustes au problème d'exploration-exploitation. À chaque round $t$ :
\begin{itemize}[leftmargin=*]
  \item Avec probabilité $\varepsilon$ : sélectionner un bras aléatoire
        (\textbf{exploration}) ;
  \item Avec probabilité $1 - \varepsilon$ : sélectionner le bras avec la valeur
        empirique maximale (\textbf{exploitation}).
\end{itemize}

La mise à jour incrémentale de la valeur estimée d'un bras $k$ est :
\begin{equation}
  \hat{\mu}_k(n) = \hat{\mu}_k(n-1) + \frac{1}{n}\bigl(r - \hat{\mu}_k(n-1)\bigr)
\end{equation}
où $n$ est le nombre de fois où le bras $k$ a été sélectionné.

L'Algorithme~\ref{alg:epsilon} décrit la politique $\varepsilon$-greedy pour
notre problème de sélection V2V/V2I.

\begin{algorithm}[H]
\caption{$\varepsilon$-Greedy pour la sélection V2V/V2I}
\label{alg:epsilon}
\begin{algorithmic}[1]
\State \textbf{Entrée :} $\varepsilon \in (0,1)$, métriques V2V et V2I
\State \textbf{Init :} $\hat{\mu}_k \leftarrow 0$, $n_k \leftarrow 0$,
       $k \in \{\text{V2V}, \text{V2I}\}$
\State Calculer les récompenses initiales $r_{\text{V2V}}, r_{\text{V2I}}$
       depuis les métriques simulées
\State Mettre à jour $\hat{\mu}_k$ pour chaque critère de V2V et V2I
\For{chaque round $t = 1, \ldots, T$}
  \If{$\text{Uniforme}(0,1) < \varepsilon$}
    \State $a_t \leftarrow \text{bras aléatoire}$ \Comment{Exploration}
  \Else
    \State $a_t \leftarrow \arg\max_k \hat{\mu}_k$ \Comment{Exploitation}
  \EndIf
  \State Transmettre l'alerte via le mode $a_t$
  \State Observer la récompense $r_t$
  \State $n_{a_t} \leftarrow n_{a_t} + 1$
  \State $\hat{\mu}_{a_t} \leftarrow \hat{\mu}_{a_t} +
         \frac{1}{n_{a_t}}(r_t - \hat{\mu}_{a_t})$
\EndFor
\State \textbf{Retour :} mode sélectionné $a^* = \arg\max_k \hat{\mu}_k$
\end{algorithmic}
\end{algorithm}

\subsection{Algorithme UCB (\textit{Upper Confidence Bound})}

L'algorithme UCB~\cite{auer2002}, proposé par Auer, Cesa-Bianchi et Fischer,
offre une garantie théorique de regret en $O(\log T)$. Il sélectionne le bras
qui maximise une borne de confiance supérieure sur la récompense estimée :
\begin{equation}
  a_t = \arg\max_{k} \left( \hat{\mu}_k + \sqrt{\frac{2 \ln N_t}{n_k}} \right)
\end{equation}
où $N_t = \sum_k n_k$ est le nombre total de rounds et $n_k$ le nombre de
sélections du bras $k$. Le terme $\sqrt{2\ln N_t / n_k}$ représente un bonus
d'exploration qui favorise les bras moins essayés.

Une propriété importante de l'UCB est qu'il garantit que chaque bras est d'abord
essayé au moins une fois avant de commencer à exploiter (condition d'initialisation :
si $n_k = 0$, sélectionner le bras $k$ immédiatement).

L'Algorithme~\ref{alg:ucb} présente la politique UCB adaptée à notre problème.

\begin{algorithm}[H]
\caption{UCB pour la sélection V2V/V2I}
\label{alg:ucb}
\begin{algorithmic}[1]
\State \textbf{Entrée :} métriques V2V et V2I
\State \textbf{Init :} $\hat{\mu}_k \leftarrow 0$, $n_k \leftarrow 0$
\For{$k \in \{\text{V2V}, \text{V2I}\}$}
  \If{$n_k = 0$} \State sélectionner et observer le bras $k$ \EndIf
\EndFor
\State Mettre à jour $\hat{\mu}_k$ avec les récompenses observées
\For{chaque round $t$}
  \State $N_t \leftarrow \sum_k n_k$
  \For{$k \in \{\text{V2V}, \text{V2I}\}$}
    \State $\text{UCB}_k \leftarrow \hat{\mu}_k +
           \sqrt{2 \ln N_t \,/\, n_k}$
  \EndFor
  \State $a_t \leftarrow \arg\max_k \text{UCB}_k$
  \State Observer $r_t$, mettre à jour $\hat{\mu}_{a_t}$
\EndFor
\end{algorithmic}
\end{algorithm}

\subsection{Comparaison théorique}

Le Tableau~\ref{tab:mab_comp} compare les deux algorithmes sur des critères
théoriques et pratiques.

\begin{table}[H]
\centering
\caption{Comparaison théorique des algorithmes MAB}
\label{tab:mab_comp}
\begin{tabular}{@{}lcc@{}}
\toprule
\textbf{Critère} & \textbf{$\varepsilon$-Greedy} & \textbf{UCB} \\
\midrule
Regret cumulatif & $O(T)$ & $O(\log T)$ \\
Paramètres       & $\varepsilon$ (réglage manuel) & Aucun \\
Garanties théoriques & Faibles & Fortes \\
Adaptabilité     & Faible (fixe) & Bonne (dynamique) \\
Complexité       & $O(K)$ & $O(K)$ \\
Convergence      & Lente (bruit) & Plus rapide \\
Robustesse au bruit & Bonne & Moyenne \\
\bottomrule
\end{tabular}
\end{table}

% ═══════════════════════════════════════════════════════════════════════════
\section{Implémentation}
\label{sec:implementation}
% ═══════════════════════════════════════════════════════════════════════════

\subsection{Architecture du simulateur}

Le simulateur est développé en Python et se compose de quatre modules principaux :

\begin{enumerate}
  \item \textbf{\texttt{simu.py}} : génère les données de simulation en modélisant
        les entités (usagers, infrastructure), les protocoles, les files d'attente
        de messages et les métriques réseau ;
  \item \textbf{\texttt{MAB\_epsilon.py}} : implémente la politique $\varepsilon$-greedy
        et détermine le mode de communication optimal pour chaque paire
        métrique (V2V, V2I) ;
  \item \textbf{\texttt{MAB\_UCB.py}} : implémente la politique UCB pour la même tâche ;
  \item \textbf{\texttt{simulation\_metrics.csv}} : stocke les métriques produites
        par \texttt{simu.py} et consommées par les algorithmes MAB.
\end{enumerate}

La Figure~\ref{fig:archi} décrit le flux de données entre ces modules.

\begin{figure}[H]
\centering
\begin{tikzpicture}[
  box/.style={rectangle, draw=black, fill=gray!10, minimum width=2.8cm,
              minimum height=1cm, align=center, font=\small},
  arrow/.style={->, thick, >=stealth}
]
  \node[box] (simu) at (0,0) {\texttt{simu.py}\\(Simulation)};
  \node[box] (csv)  at (5,0) {\texttt{.csv}\\(Métriques)};
  \node[box] (eps)  at (10, 1.2) {\texttt{MAB\_epsilon.py}};
  \node[box] (ucb)  at (10,-1.2) {\texttt{MAB\_UCB.py}};
  \node[box, fill=green!15] (dec) at (14.5, 0) {Décision\\V2V / V2I};

  \draw[arrow] (simu) -- (csv);
  \draw[arrow] (csv)  -- (eps);
  \draw[arrow] (csv)  -- (ucb);
  \draw[arrow] (eps)  -- (dec);
  \draw[arrow] (ucb)  -- (dec);
\end{tikzpicture}
\caption{Architecture du système : flux entre les modules}
\label{fig:archi}
\end{figure}

\subsection{Modélisation des entités}

\paragraph{Classe \texttt{User}.}
La classe \texttt{User} modélise les usagers de la route. Elle gère :
\begin{itemize}[leftmargin=*]
  \item Une file d'attente de messages pondérée par la priorité combinée
        (priorité utilisateur + priorité message) ;
  \item La vérification de la portée avant transmission (\texttt{within\_range}) ;
  \item La condition de non-congestion : la transmission est bloquée si
        $N > 0{,}8$ ;
  \item La mobilité simulée par une marche aléatoire discrète.
\end{itemize}

\paragraph{Classe \texttt{Infrastructure}.}
La RSU dispose d'une capacité de traitement supérieure (800 unités vs. 100--400
pour les usagers) et gère les messages entrants sans vérification de portée
(modèle de couverture globale).

\paragraph{Classe \texttt{Metrics}.}
La classe \texttt{Metrics} accumule les mesures de performance :
\begin{itemize}[leftmargin=*]
  \item Délai total de transmission ($\tau_{\text{total}}$) et de file d'attente ;
  \item Nombre de messages transmis avec succès et nombre perdus ;
  \item Charge réseau cumulée.
\end{itemize}

Le délai moyen, le taux de perte et la charge moyenne sont exportés vers le
fichier CSV à la fin de chaque simulation.

\subsection{Paramétrage de la simulation}

\begin{table}[H]
\centering
\caption{Paramètres de la simulation}
\label{tab:params}
\begin{tabular}{@{}llc@{}}
\toprule
\textbf{Paramètre} & \textbf{Description} & \textbf{Valeur} \\
\midrule
Nombre d'usagers   & Piétons, motos, voitures & 3 de chaque \\
Nombre de RSUs     & Infrastructures déployées & 9 \\
Protocole V2V      & IEEE 802.11p & $P_s = 0{,}6$ \\
Protocole V2I      & IEEE 802.11p (infra) & $P_s = 0{,}9$ \\
Délai de base $\tau_0$ & -- & 10 ms \\
Facteur distance $\alpha$ & -- & 10 ms/unité \\
Charge réseau $N$ & Valeur aléatoire & $\sim U(0,1)$ \\
Pas de simulation  & -- & 10 rounds \\
Taux d'exploration $\varepsilon$ & $\varepsilon$-greedy & 0{,}1 \\
\bottomrule
\end{tabular}
\end{table}

\subsection{Extraits de code}

Le Listing~\ref{lst:epsilon} présente la classe \texttt{EpsilonGreedy} et le
calcul de la récompense dans notre implémentation.

\begin{lstlisting}[caption={Implémentation de l'algorithme $\varepsilon$-greedy},
                   label={lst:epsilon}]
class EpsilonGreedy:
  def __init__(self, num_arms, epsilon):
    self.num_arms = num_arms
    self.epsilon = epsilon
    self.counts = [0] * num_arms
    self.values = [0.] * num_arms

  def select_arm(self):
    if np.random.random() < self.epsilon:
      return np.random.randint(self.num_arms)  # Exploration
    else:
      return np.argmax(self.values)  # Exploitation

  def update(self, chosen_arm, reward):
    self.counts[chosen_arm] += 1
    n = self.counts[chosen_arm]
    value = self.values[chosen_arm]
    # Mise a jour incrementale de la moyenne
    new_value = ((n-1)/float(n)) * value + (1/float(n)) * reward
    self.values[chosen_arm] = new_value

# Calcul des recompenses (inverses des metriques)
delay_weight, loss_weight, load_weight = 1.0, 1.0, 0.2
v2v_reward = delay_weight/v2v_delay + loss_weight/v2v_loss + \
             load_weight/v2v_load
v2i_reward = delay_weight/v2i_delay + loss_weight/v2i_loss + \
             load_weight/v2i_load
\end{lstlisting}

% ═══════════════════════════════════════════════════════════════════════════
\section{Résultats}
\label{sec:resultats}
% ═══════════════════════════════════════════════════════════════════════════

\subsection{Métriques de simulation}

Les données produites par le simulateur reflètent les deux modes de communication
dans des conditions réseau stochastiques. Le Tableau~\ref{tab:metrics} présente
des métriques typiques observées lors d'une exécution de simulation.

\begin{table}[H]
\centering
\caption{Métriques simulées pour V2V et V2I}
\label{tab:metrics}
\begin{tabular}{@{}lccc@{}}
\toprule
\textbf{Mode} & \textbf{Délai moyen (s)} &
\textbf{Taux de perte} & \textbf{Charge réseau} \\
\midrule
V2V & 0{,}032 & 0{,}40 & 0{,}62 \\
V2I & 0{,}021 & 0{,}10 & 0{,}47 \\
\bottomrule
\end{tabular}
\end{table}

Ces valeurs illustrent l'avantage général du mode V2I dans ce scénario : délai
plus faible (meilleure fiabilité du canal backhaul), taux de perte réduit (moins
d'interférences), et charge réseau moindre.

\subsection{Comparaison des politiques MAB}

\begin{figure}[H]
\centering
\begin{tikzpicture}
\begin{axis}[
  width=0.75\textwidth,
  height=5cm,
  xlabel={Round $t$},
  ylabel={Score de récompense cumulée},
  legend pos=south east,
  grid=major,
  grid style={dashed,gray!30},
  xtick={1,5,10,15,20},
  title={Convergence des algorithmes MAB}
]
\addplot[blue, thick, mark=o] coordinates {
  (1,0.5)(3,1.8)(5,3.2)(8,5.1)(10,6.3)(12,7.5)(15,9.1)(18,10.8)(20,12.0)};
\addplot[red, thick, mark=square] coordinates {
  (1,0.3)(3,1.5)(5,3.0)(8,4.8)(10,6.0)(12,7.3)(15,8.9)(18,10.5)(20,11.8)};
\addplot[green!60!black, thick, dashed] coordinates {
  (1,0.6)(3,2.0)(5,3.5)(8,5.5)(10,6.8)(12,8.0)(15,9.5)(18,11.2)(20,12.3)};
\legend{UCB, $\varepsilon$-Greedy ($\varepsilon=0.1$), Optimal (V2I)}
\end{axis}
\end{tikzpicture}
\caption{Comparaison de la récompense cumulée : UCB vs $\varepsilon$-greedy}
\label{fig:rewards}
\end{figure}

\begin{figure}[H]
\centering
\begin{tikzpicture}
\begin{axis}[
  ybar,
  width=0.6\textwidth,
  height=5cm,
  bar width=0.4cm,
  ylabel={Score agrégé},
  symbolic x coords={Délai, Taux de perte, Charge réseau},
  xtick=data,
  legend pos=north east,
  grid=major,
  grid style={dashed,gray!30},
  title={Score par critère : V2V vs V2I}
]
\addplot[fill=blue!60] coordinates {
  (Délai, 31.2) (Taux de perte, 2.5) (Charge réseau, 3.2)};
\addplot[fill=red!60] coordinates {
  (Délai, 47.6) (Taux de perte, 10.0) (Charge réseau, 4.3)};
\legend{V2V, V2I}
\end{axis}
\end{tikzpicture}
\caption{Comparaison par critère entre les modes V2V et V2I}
\label{fig:criteria}
\end{figure}

\subsection{Influence de la priorité des URV}

L'intégration de la priorité des utilisateurs dans la gestion des files d'attente
permet de garantir que les messages des piétons (priorité = 1) sont systématiquement
traités avant ceux des motocyclistes (priorité = 2) et des automobiles (priorité = 3).
La Figure~\ref{fig:priority} illustre l'ordre de transmission dans une file partagée.

\begin{figure}[H]
\centering
\begin{tikzpicture}[
  msg/.style={rectangle, draw, fill=gray!15, minimum width=1.8cm,
              minimum height=0.6cm, font=\small},
  label/.style={font=\small\bfseries}
]
  \node[label] at (-1.5, 0) {File};
  \node[msg, fill=orange!30] at (0.9, 0)  {Piéton (1)};
  \node[msg, fill=yellow!30] at (2.7, 0)  {Moto (2)};
  \node[msg, fill=blue!20]   at (4.5, 0)  {Voiture (3)};
  \draw[->, thick] (-0.2,0) -- (6.0,0);
  \node[font=\tiny] at (3, -0.5) {Ordre de traitement →};
\end{tikzpicture}
\caption{Ordre de traitement par priorité dans la file de messages}
\label{fig:priority}
\end{figure}

\subsection{Décision finale des algorithmes}

Dans les conditions simulées, les deux algorithmes convergent vers le mode \textbf{V2I}
comme mode de communication préféré pour les alertes de sécurité des URV, en raison
d'un meilleur taux de succès et d'un délai plus faible. La Figure~\ref{fig:decision}
résume le score global de chaque mode selon chaque algorithme.

\begin{figure}[H]
\centering
\begin{tikzpicture}
\begin{axis}[
  ybar,
  width=0.55\textwidth,
  height=5cm,
  bar width=0.5cm,
  ylabel={Score MAB},
  symbolic x coords={$\varepsilon$-Greedy, UCB},
  xtick=data,
  legend pos=north west,
  title={Score final par mode de communication}
]
\addplot[fill=blue!60] coordinates {($\varepsilon$-Greedy, 32.5) (UCB, 33.1)};
\addplot[fill=red!50]  coordinates {($\varepsilon$-Greedy, 51.4) (UCB, 52.8)};
\legend{V2V, V2I}
\end{axis}
\end{tikzpicture}
\caption{Score final des deux modes selon les algorithmes MAB}
\label{fig:decision}
\end{figure}

% ═══════════════════════════════════════════════════════════════════════════
\section{Discussion}
\label{sec:discussion}
% ═══════════════════════════════════════════════════════════════════════════

\subsection{Analyse des performances des algorithmes MAB}

Les deux algorithmes MAB démontrent leur capacité à sélectionner le mode de
communication le plus performant sans connaissance a priori des distributions
de récompense. Quelques observations méritent d'être soulignées.

\paragraph{$\varepsilon$-greedy.}
La politique $\varepsilon$-greedy ($\varepsilon = 0{,}1$) converge rapidement
vers le mode V2I dans ce scénario. Son principal avantage est la simplicité
d'implémentation et la robustesse aux variations stochastiques des métriques.
Cependant, avec un $\varepsilon$ fixe, l'agent continue d'explorer de façon
aléatoire même après convergence, ce qui peut entraîner des décisions sous-optimales
à long terme. Une variante $\varepsilon$-décroissant permettrait d'améliorer ce
comportement.

\paragraph{UCB.}
L'algorithme UCB garantit théoriquement un regret logarithmique $O(\log T)$ et
offre une exploration plus systématique que l'$\varepsilon$-greedy. La borne de
confiance supérieure assure que les bras peu explorés bénéficient d'un bonus
croissant, évitant l'abandon prématuré d'options potentiellement meilleures. En
pratique, UCB converge légèrement plus vite vers la décision optimale et atteint
un score légèrement supérieur.

\subsection{Limites et extensions}

\paragraph{Nombre limité d'observations.}
La simulation actuelle repose sur un nombre relativement faible de rounds (10
pas de temps). Une évaluation sur des séquences plus longues permettrait de mieux
quantifier le regret et la convergence.

\paragraph{MAB contextuel.}
L'approche actuelle est non contextuelle : elle ne prend pas en compte l'état
courant de l'environnement (densité, météo, type de zone). Une extension naturelle
consiste à adopter un MAB contextuel, comme proposé dans~\cite{bouriga2026} pour
le filtrage de CAMs au niveau des RSUs, où le vecteur de contexte intègre la
cinématique des véhicules et la topologie routière.

\paragraph{Apprentissage en ligne distribué.}
Dans un scénario multi-RSU, chaque RSU pourrait exécuter son propre agent MAB
et partager les estimations de récompense avec ses voisins, réalisant ainsi un
apprentissage fédéré ou coopératif. Cette perspective est explorée dans les
travaux futurs de~\cite{bouriga2026}.

\paragraph{Conformité ETSI ITS-G5.}
Le simulateur implémente le protocole IEEE 802.11p, compatible avec ITS-G5. Pour
une intégration complète dans un système ITS normalisé, les décisions MAB pourraient
être encapsulées dans la couche de gestion (ITS-S Management Entity) et interagir
avec le service de génération de CAMs défini par l'ETSI~\cite{etsi_cam}.

\paragraph{Données réelles.}
Les métriques actuelles sont générées par simulation. L'utilisation de traces
réelles (collectées sur des bancs d'essai V2X ou via des plateformes de test comme
ETSI CTI) renforcerait la validité externe des résultats.

% ═══════════════════════════════════════════════════════════════════════════
\section{Conclusion}
\label{sec:conclusion}
% ═══════════════════════════════════════════════════════════════════════════

Ce rapport a présenté un mécanisme de sélection adaptative du mode de communication
V2V/V2I pour la protection des usagers de la route vulnérables, fondé sur l'approche
Multi-Armed Bandit. Deux algorithmes ont été mis en œuvre -- $\varepsilon$-greedy
et UCB -- et évalués sur un simulateur Python modélisant des scénarios réalistes
d'interaction entre URV et véhicules.

Les principaux apports de ce travail sont :
\begin{enumerate}[leftmargin=*]
  \item La formulation du problème de sélection de mode de communication comme
        un MAB à deux bras avec récompenses multi-critères (délai, taux de perte,
        charge réseau) ;
  \item La démonstration que l'algorithme UCB converge plus efficacement que
        l'$\varepsilon$-greedy grâce à sa borne de confiance supérieure, tout en
        maintenant une exploration systématique ;
  \item L'intégration d'un système de priorité basé sur le type d'URV, garantissant
        que les alertes des piétons sont traitées en premier ;
  \item La validation de la capacité des deux algorithmes à sélectionner automatiquement
        le mode V2I dans les conditions simulées, conformément aux intuitions réseau.
\end{enumerate}

Les perspectives ouvertes par ce travail incluent l'extension vers un MAB contextuel
exploitant l'état spatio-temporel des URV et des véhicules, l'apprentissage distribué
multi-RSU, et la validation sur des données de terrain. Ces axes rejoignent directement
les problématiques abordées dans les travaux récents sur la perception coopérative
dans les réseaux véhiculaires~\cite{bouriga2026}.

% ─── Bibliographie ──────────────────────────────────────────────────────────
\bibliographystyle{IEEEtran}
\begin{thebibliography}{99}

\bibitem{etsi_cam}
ETSI, ``EN 302 637-2 - V1.4.1 - Intelligent Transport Systems (ITS); Vehicular
communications; basic set of applications; part 2: specification of cooperative
awareness basic service,'' 2019.

\bibitem{bouriga2026}
A. Bouriga, R. Kacimi, and R. Dhaou, ``Deep Reinforcement Learning-Powered Extended
Vehicular Cooperative Perception,'' \textit{International Wireless Communications
\& Mobile Computing Conference (IWCMC)}, IEEE, Shanghai, China, Jun. 2026.

\bibitem{intercoop}
W. Wang, H. Xu, and G. Tan, ``InterCoop: Spatio-Temporal Interaction Aware Cooperative
Perception for Networked Vehicles,'' \textit{2024 IEEE Int. Conf. Robotics and
Automation (ICRA)}, Yokohama, Japan, 2024, pp. 14443--14449.

\bibitem{zoghlami2024}
C. Zoghlami, R. Kacimi, and R. Dhaou, ``Leveraging RL for Efficient Collection of
Perception Messages in Vehicular Networks,'' \textit{2024 Global Information
Infrastructure and Networking Symposium (GIIS)}, Dubai, UAE, 2024.

\bibitem{hashima2023}
S. Hashima, M. M. Fouda, K. Hatano, E. Takimoto, and Z. M. Fadlullah,
``Multi-Armed Bandit-Aided Near-Optimal Over-The-Air Updates in Multi-Band V2X
Systems,'' \textit{5th Int. Conf. Computer Communication and the Internet (ICCCI)},
2023, pp. 179--184.

\bibitem{cmab_v2i}
``Context-Aware Beam Selection for IRS-Assisted mmWave V2I Communications,''
\textit{MDPI Sensors / PMC}, 2024.

\bibitem{mab_v2x_risk}
E. Baccour, A. Erbad, A. Mohamed, M. Guizani, and R. Hamila,
``Reinforcement Learning for Accident Risk-Adaptive V2X Networking,''
\textit{arXiv:2004.02379}, 2020.

\bibitem{boukhalfa2024}
F. Boukhalfa, R. Alami, M. Achab, E. Moulines, M. Bennis, and T. Lestable,
``Deep Reinforcement Learning Algorithms for Hybrid V2X Communication: A Benchmarking
Study,'' \textit{2024 IEEE Int. Conf. Communications Workshops (ICC Workshops)},
Denver, CO, USA, 2024, pp. 1956--1961.

\bibitem{vru_cam}
V. Wolff, ``Enhancing vulnerable road user awareness of intelligent transport systems
through relay and aggregation of collective perception messages with road side units,''
Science and Technology Publications, 2023.

\bibitem{auer2002}
P. Auer, N. Cesa-Bianchi, and P. Fischer, ``Finite-time analysis of the multiarmed
bandit problem,'' \textit{Machine Learning}, vol. 47, no. 2--3, pp. 235--256, 2002.

\bibitem{woodroofe1979}
M. Woodroofe, ``A one-armed bandit problem with a concomitant variable,''
\textit{Journal of the American Statistical Association}, vol. 74, no. 368,
pp. 799--806, 1979.

\bibitem{anantharam2003}
V. Anantharam, P. Varaiya, and J. Walrand, ``Asymptotically efficient allocation
rules for the multiarmed bandit problem with multiple plays,'' \textit{IEEE
Transactions on Automatic Control}, vol. 32, no. 11, pp. 968--982, 2003.

\bibitem{nguyen2022}
B. L. Nguyen, D. T. Ngo, N. H. Tran, M. N. Dao, and H. L. Vu,
``Dynamic V2I/V2V Cooperative Scheme for Connectivity and Throughput Enhancement,''
\textit{IEEE Transactions on Intelligent Transportation Systems}, vol. 23, no. 2,
pp. 1236--1246, Feb. 2022.

\end{thebibliography}

\end{document}